% PACKAGES AND OTHER DOCUMENT CONFIGURATIONS

\documentclass{resume} % Use the custom resume.cls style

\usepackage{enumitem}
\usepackage{gensymb}
\usepackage[left=0.4 in,top=0.4in,right=0.4 in,bottom=0.4in]{geometry} % Document margins
\usepackage{xeCJK} % 中文支持（XeLaTeX）
\setCJKmainfont{PingFang SC}[BoldFont=PingFang SC, BoldFeatures={FakeBold=2.5}]
\usepackage{url}
\newcommand{\link}[1]{\url{#1}}
\newcommand{\tab}[1]{\hspace{.2667\textwidth}\rlap{#1}} 
\newcommand{\itab}[1]{\hspace{0em}\rlap{#1}}
\renewcommand{\baselinestretch}{1}
\name{CHEN XINGQIANG} % Your name
\address{杭州，中国} % Your secondary addess (optional)
\address{+86 150-0133-4125 \\ joy6677@outlook.com \\ github.com/chenxingqiang}

\begin{document}

% OBJECTIVE

\begin{rSection}{求职意向}
{寻求{\bf 全职}职位：
{\bf AI 加速器编译器 \/ 系统软件} \/
{\bf 大语言模型推理系统} \/
{\bf 软硬件协同设计}。}
\end{rSection}

\begin{rSection}{交叉学科背景}
{理论物理与生物物理训练，贯通计算化学、分子动力学与计算机科学。竞赛：{\bf 全国大学生数学建模竞赛国家一等奖}（2012）、{\bf 北美大学生数学建模竞赛（MCM）一等奖}（2013）。}
\end{rSection}

% EDUCATION SECTION
%---------------------------------------------------------------------

\begin{rSection}{教育背景}

{\bf 生物物理学与凝聚态物理博士} \hfill {2014年9月 - 2018年7月 }\\ 
厦门大学，物理与机电工程学院 GPA : 3.5/4.0
\\
\\
{\bf 理论物理学学士} \hfill {2010年9月 - 2014年7月}\\
厦门大学，物理与机电工程学院 GPA : 3.0/4.0

\end{rSection}


% TECHNICAL STRENGTHS SECTION
%---------------------------------------------------------------------

\begin{rSection}{技能与项目}

 \item {\bf 编程语言}：精通 C++、Python，熟悉 Java、Scala、Go、CUDA。

 \item {\bf MLOps工具}：熟练使用 PyTorch、TensorFlow、Triton、TVM、Kubernetes；Hadoop、Flink、Spark 等。
 
 \item {\bf 大语言模型 \/ 推理系统}：Transformers、KV-Cache、量化、批处理与调度、Triton内核、CUDA、FlashAttention、图融合、GPU Profiling 与 Replay。

 \item {\bf AI编译器 \/ 内核}：MLIR 图编译、CIM/NPU CodeGen、异构运行时；{\bf Palladium} 仿真验证与 HW/SW 协同 Debug；YiRage YPK；TraceSmith。

 \item {\bf 深度学习模型}：熟悉 Transformers、RL、GNN、DNN、CNN 等。

\item {\bf 跨学科专长}：具有物理学、计算生物学、计算化学、分子动力学、量子力学仿真和计算机科学背景。

\item{\bf 开源亮点}：

\begin{itemize}
  \item {\bf TraceSmith} — 面向 AI 编译器与深度学习框架的跨平台 GPU Profiling \& Replay 系统（CUDA/ROCm/Metal/MACA/Ascend；无锁采集、确定性回放；PyPI 发布）。
  
  \link {https://github.com/chenxingqiang/TraceSmith}
  
  \item {\bf Research CLI} — 具有学术研究功能的命令行 AI 工作流工具；支持多模型提供商和 MCP 工具。
  
  \link {https://github.com/iechor-research/research-cli}
  
  \item {\bf JAX-sklearn} — 具有自动 JAX 加速功能的 scikit-learn 替代方案（{\textit{13{,}058 项测试通过；PyPI 发布}}）。 
  
  \link {https://github.com/chenxingqiang/jax-sklearn}
  
  \item {\bf YouTube-DNN-RecSys} — YouTube 风格推荐模型的简洁 TensorFlow 与 Spark 实现（Python+Scala）。 
  
  \link {https://github.com/chenxingqiang/YouTube-DNN-RecSys}
  
  \item {\bf Questions-Gen} — 数学竞赛题目生成器（ Qwen3-14B、RL GRPO、KD；CLI 与 HF 模型）。 
  
  \link {https://github.com/chenxingqiang/GenQuestion}

  \item {\bf DeepMicroPath} — 多模式研究 Agent，提供微生物学专业模式；以非侵入式包装器实现 auto/default/microbiology Prompt，支持 Docker/SLURM 部署。
  
  \link {https://github.com/chenxingqiang/DeepMicroPath}
  
  \item {\bf Medical-NER-Qwen-4B-Thinking} — 通过LoRA微调的Qwen3-4B-Thinking用于医疗实体与关系抽取（英语）。
  
  \link {https://huggingface.co/xingqiang/Medical-NER-Qwen-4B-Thinking}
  
  \item {\bf QuestionsGen-Qwen3-14B-Final (FP)} — 使用Unsloth训练的全精度合并最终模型；部署就绪的HF权重。 
  
  \link {https://huggingface.co/xingqiang/QuestionsGen-Qwen3-14b-final-fp-merged}
  
  \item {\bf YiRage持久化内核(YPK)} — LLM超级内核编译器。
  
  \link {https://github.com/chenxingqiang/YiRage}
\end{itemize}

\end{rSection}


\begin{rSection}{工作经历}

\begin{rSubsection}{编译器\&工具链专家，知存算力（Witmem）}{杭州，中国}
{\textbf{编译器 / Palladium 验证}}{2025年12月10日 - 至今}
\item {负责存算芯片产品线编译器与 Palladium 仿真验证，覆盖模型编译、指令生成、芯片执行及 Emulation 缺陷定位。}\\
\item {完成 Palladium 软件 Bring-up：Compiler 输出在仿真器上运行，结合 GDB 在仿真模型、芯片与系统软件层做联合调试。}\\
\item {实现{\textbf{CIM-MLIR}}异构编译器：PyTorch/TensorFlow/ONNX 经 TOSA/StableHLO 映射至 CPU、GPU 与 CIM，含前端、图划分、硬件后端与运行时。}\\
\item {实现代价驱动图划分与 CIMHigh/Mid/Low 三级 IR；用 OR-Tools 分配 Die/L2/L0 资源，生成可烧录{\textbf{config.bin}}、权重{\textbf{map.bin}}，以及 GPU PTX 与 LLVM 主机代码。}\\
\item {实现跨设备运行时：分设备执行器、统一内存、计算与传输重叠及三级缓冲；提供 C++/Python 接口及延迟、利用率、传输开销统计。}\\
\item {实现{\textbf{GRC-PIM}}：PIM 上基于 drift gating 与 GEMV 复用的残差校准及软判决纠错，投稿 {\textit{npj Unconventional Computing}}。}\\
\end{rSubsection}
\vspace{0.5em}

\begin{rSubsection}{高级AI编译器工程师，亿铸智能}{杭州，中国}
{\textbf{研发部门}}{2022年9月 - 2025年12月}
\item {主导开发{\textbf{YiRage 持久化内核(YPK)}}，将端到端LLM推理编译成单个超级内核，在Llama2-7B上实现{\textbf{1.2×–6.7×延迟缩减}}，GPU利用率提升{\textbf{40\%}}，显存占用减少{\textbf{25\%}}。}\\
\item {设计并实现融合算子架构（RMSNorm + Linear 等），消除{\textbf{85\%}}的内核启动开销，单次推理延迟从{\textbf{120ms 降至 18ms}}（批大小=1）。}\\
\item {构建自适应调度系统，整合 Triton 与 FlashInfer 内核，支持多GPU扩展至{\textbf{8卡并行}}，吞吐量达到{\textbf{15,000 tokens/秒}}。}\\
\item {开发图编译优化器，实现自动内核参数调优和KV-cache内存布局优化，内存访问效率提升{\textbf{60\%}}，缓存命中率达到{\textbf{92\%}}。}\\
\item {技术成果转化：产品已部署到{\textbf{2个}}数据中心客户。}\\
\item {在 ReRAM 芯片公司主要负责编译器后端适配与本公司软件生态构建，聚焦公司 SoC 的 LLVM 后端 Target 研发，推进代码规范化与复杂指令集架构开发。}\\
\item {承担 AI 模型（含 LLMs）的转换与细粒度分析，适配多种运行时格式，评估推理阶段算子指标（Cycles/MACs 等）与计算效率；集成 TVM、ONNX、MLCommons、TensorFlow/PyTorch 于 YZlite，对多种模型运行时进行算子级指标统计并输出各类芯片上的 benchmark。}\\
\item {YiRage 持久化内核 (YPK) — LLM 超级内核编译器与运行时，实现 1.2×–6.7× 延迟缩减。}\\
\link{https://github.com/chenxingqiang/YiRage}\\
\item {ReRAM-YZLite — 亿铸科技 ReRAM 芯片的机器学习模型开发工具包，提供命令行接口、Python API、模型分析器、训练评估、量化等功能，支持嵌入式平台的 ML 模型开发。}\\
\link{https://github.com/ReRAM-Labs/yzlite}\\
\link{https://reram-labs.github.io/yzlite}\\
\end{rSubsection}
\vspace{0.5em}

\begin{rSubsection}{隐私保护算法工程师，蓝象智联}{杭州，中国}
{\textbf{IT Group}}{2022年1月 - 2022年9月}
\item {构建基于同态加密的联邦学习平台，支持{\textbf{20+}}金融机构数据协作，在保护隐私前提下模型精度提升{\textbf{15\%}}，训练时间缩短{\textbf{40\%}}。}\\
\item {设计多方安全计算协议，为{\textbf{8家}}银行提供风控模型联合训练服务，欺诈检测准确率达到{\textbf{96.7\%}}，误报率降低{\textbf{30\%}}。}\\
\item {开源贡献：开发{\textbf{Secret-Learn}}隐私保护机器学习库，实现{\textbf{573个}}算法（191算法 × FL/SS/SL三种隐私模式），100\% sklearn API兼容，集成JAX加速（{\textbf{5×+}}性能提升）和SecretFlow隐私计算框架，支持联邦学习、秘密分享、分割学习等多种隐私保护场景（{\textbf{PyPI发布，150,000+行代码}}）。}\\
\link{https://github.com/chenxingqiang/Secret-Learn}\\
\item {技术创新产出：申请{\textbf{3项}}核心算法专利，核心技术被{\textbf{5个}}商业产品采用，创造营收{\textbf{800万}}人民币。}\\
\item {基于秘密分享的纵向联邦线性支持向量分类方法（CN114925786A）。}\\
\link{https://patents.google.com/patent/CN114925786A/en}\\
\item {随机森林训练方法及装置（CN114997345A）。}\\
\link{https://patents.google.com/patent/CN114997345A/en}\\
\item {用于联邦学习的数据处理方法及装置（CN115130681A）。}\\
\link{https://patents.google.com/patent/CN115130681A/en}\\
\end{rSubsection}
\vspace{0.5em}
\begin{rSubsection}{T7 算法工程专家，百图生科}{北京，中国}
{\textbf{算法工程组}}{2021年9月 - 2021年12月}
\item {疫情期间加入AI制药公司，负责构建生物信息学工具链平台，整合多种生物信息学算法和工具，支持蛋白质序列分析、结构预测、功能注释等核心流程。}\\
\item {设计并实现生物算法数据平台，建立统一的数据管理和处理框架，支持大规模生物数据的存储、检索和分析，提升算法研发效率。}\\
\item {开发蛋白质结构预测模型分析系统，集成AlphaFold等先进模型，构建端到端的蛋白质结构预测与验证流程，为药物设计提供结构基础。}\\
\item {技术产出：开源AlphaFold2代码库与学习资源（{\textbf{26 stars}}），包含源代码分析、算法实现、应用案例和论文分享PPT，为AlphaFold2研究提供完整参考。}\\
\link{https://github.com/chenxingqiang/alphafold2-codecs}\\
\item {获得{\textbf{2项}}软件著作权：}\\
\item {百图生科生信工具自动化平台（登记号：2021SR2004464）。}\\
\item {Biomap生物医药信息平台（登记号：2021SR1862397）。}\\
\end{rSubsection}
\vspace{0.5em}
\begin{rSubsection}{高级算法工程师，Yalla 集团(纽交所: YALA)}{杭州，中国}
{\textbf{数据中心工程组}}{2020年9月 - 2021年9月}
\item {重构Yalla聊天室推荐系统，采用多目标优化算法，用户日活跃时长增加{\textbf{35\%}}，7日留存率从{\textbf{42\%}}提升至{\textbf{61\%}}，月收入增长{\textbf{1200万}}美元。}\\
\item {构建Ludo智能风控系统，结合实时行为分析和机器学习模型，恶意用户识别准确率达{\textbf{94\%}}，退款率降低{\textbf{45\%}}，挽回损失约{\textbf{500万}}美元。}\\
\item {开发德州扑克AI机器人，基于深度强化学习，胜率达到{\textbf{78\%}}，驱动用户游戏时长增加{\textbf{25\%}}，付费转化率提升{\textbf{18\%}}。}\\
\end{rSubsection}
\vspace{0.5em}
\begin{rSubsection}{高级算法工程师，格家网络科技}{杭州，中国}
{\textbf{搜索推荐工程组}}{2019年3月 - 2020年9月}
\item {设计多场景推荐算法架构，整合协同过滤、深度学习和知识图谱技术，商品点击率(CTR)提升{\textbf{28\%}}，转化率(CVR)增长{\textbf{35\%}}，GMV增长{\textbf{4200万}}人民币。}\\
\item {优化实时推荐引擎，支持{\textbf{100万}}QPS并发，响应延迟控制在{\textbf{50ms}}以内，推荐覆盖率达到{\textbf{95\%}}，用户满意度评分从{\textbf{3.2}}提升至{\textbf{4.6}}。}\\
\item {YouTube-DNN-RecSys — YouTube 风格推荐模型的简洁 TensorFlow 与 Spark 实现（Python+Scala）（{\textbf{60 stars}}）。}\\
\link{https://github.com/chenxingqiang/YouTube-DNN-RecSys}\\
\end{rSubsection}
\vspace{0.5em}
\begin{rSubsection}{机器学习工程师，便利蜂科技}{北京，中国}
{\textbf{数据工程组}}{2017年3月 - 2019年3月}
\item {构建智能销售预测系统，覆盖{\textbf{50+}}门店，预测精度达到{\textbf{92\%}}，库存周转率提升{\textbf{30\%}}，减少缺货损失{\textbf{600万}}人民币，降低滞销损失{\textbf{400万}}人民币。}\\
\item {开发用户画像与商品关联分析系统，处理{\textbf{1000万+}}用户行为数据，客户身份识别准确率{\textbf{89\%}}，个性化推荐带来销售增长{\textbf{22\%}}。}\\
\item {设计计算机视觉商品识别系统，支持{\textbf{500+}}SKU自动识别，识别准确率{\textbf{96\%}}，收银效率提升{\textbf{60\%}}，顾客排队时间减少{\textbf{40\%}}。}\\
\end{rSubsection}



\end{rSection}

\begin{rSection}{研究与实习}
\begin{rSubsection}{\bf 蛋白质分子动力学仿真与QM/MM模型  }{厦门大学，中国}
{\bf (国家自然科学基金项目) }{2014年9月 - 2017年3月}
  \item {利用量子力学进行蛋白质动力学仿真研究，分析蛋白质生物信息。涉及理论物理、结构生物学、计算化学和计算机科学等跨学科知识。}\\
\end{rSubsection}

\begin{rSubsection}{\bf 液体活检ct-DNA试剂盒开发与DNA甲基化数据分析  }{乐普医疗，中国}
{\bf (研发工程师实习) }{2016年3月 - 2017年3月}
  \item {从事液体活检试剂盒开发，需要特定的癌症筛查生物标志物。利用现有基因芯片数据(TCGA\&GEO)，分析10,000个样本以定位人类Septin9基因上的甲基化位点，这是结肠癌的早期指标。这为开发更准确的结肠癌临床诊断试剂盒提供了信息。}\\
\end{rSubsection}

\end{rSection}
%----------------------------------------------------------------------------------------

\begin{rSection}{获奖情况}

\begin{rSubsection}{}{}
{}{}
  \item {评委，应用算法最佳实践初审，2021年。}\\
  \item {第57名和第32名，共1739队，IEEE ICDM全球气象AI挑战赛，2018年。}\\
  \item {一等奖，北美大学生数学建模竞赛，2013年。}\\
  \item {第6名，国际太阳能十项全能竞赛，2013年。}\\
  \item {国家一等奖，全国大学生数学建模竞赛，2012年。}\\
\end{rSubsection}

\end{rSection}

%----------------------------------------------------------------------------------------
\begin{rSection}{发表论文}
\begin{rSubsection}{}{}
{}{}
  \item {Chen X. GRC-PIM: A Drift-Gated, GEMV-Reusing Residual Calibrator for Soft-Decision Error Correction in Processing-in-Memory. {\textit{npj Unconventional Computing}}（投稿中）。}\\
  \item {Cai Y, Su M, Ahmad A, Hu X, Sang J, Kong L, Chen X, Wang C, Shuai J, Han A. Conformational dynamics of the essential sensor histidine kinase WalK. {\textit{Acta Crystallographica Section D: Structural Biology.}} 2017 Oct 1;73(10):793-803.}\\
  \item {Xia H, Chen X, Chen X, Cheng J, Liu Y, Chen X, Zhang Q, Zou G. Chiral discrimination of amino acid enantiomers based on different interactions with Cu2+. {\textit{Sensors and Actuators B: Chemical}}. 2018 Jan 1;254:44-51.}\\
  \item {Ahmad A, Cai Y, Chen X, Shuai J, Han A. Conformational dynamics of response regulator RegX3 from Mycobacterium tuberculosis. {\textit{PloS one}}. 2015 Jul 22;10(7):e0133389.}\\
  \item {Chen Y, Qi H, Li X, Cai M, Chen X, Liu W, Shuai J. Suppressing effect of C a 2+ blips on puff amplitudes by inhibiting channels to prevent recovery. {\textit{Physical Review E}}. 2016 Aug 15;94(2):022411.}\\
\end{rSubsection}
\end{rSection}


%----------------------------------------------------------------------------------------

% end CJK env removed; using xeCJK
\end{document}
